% PR-D — first-class bug taxonomy section.
%
% The 17-bug taxonomy is the central methodological contribution of
% the kaggle2 study.  Bugs 1–13 are the silent-bug guards that ship
% in the original assigner pipeline; bugs 14–17 are PR-C-era
% additions (anchor-extender warmup, priors_v3 Bahasa over-fire, KD
% pooling on 0-box receipts, RAG self-hit on val).  Each row carries
% a ΔF1 measurement from the ablation grid plus the regression-guard
% test name so a reviewer can cross-reference the test suite.
\section{Silent-bug taxonomy}\label{sec:taxonomy}

Receipt-KIE pipelines suffer from a long tail of silent bugs:
behaviour-altering defects that pass every unit test, propagate
through training, and surface only as an unexplained F1 deficit.
We catalogue every silent bug we encountered while building
\textsc{kaggle2} and pair each with (i) a measured $\Delta$F1 from
the ablation grid (see Sec.~\ref{sec:experiments}), (ii) a one-line
detection method, and (iii) a pinned regression-guard test name.

\begin{table}[t]
\caption{Silent-bug taxonomy (1--17). $\Delta$F1 is the measured drop
when the corresponding guard is disabled, averaged across
\VAR{n_trials}~seeds. Regression guard names are pytest test ids in
\texttt{tests/}; a reviewer can run any single guard via
\texttt{pytest tests/<name>.py}.}\label{tab:taxonomy}
\centering
\input{tab_bug_taxonomy}
\end{table}

The first 13 rows of Table~\ref{tab:taxonomy} are the original
silent-bug guards reported in our v3 submission; rows 14--17 are new
to this work.  Bug~14 captures an anchor-extender warmup ordering
defect that causes the address head's extender branch to diverge
from the anchor logits when training-time warmup is shorter than
the head's gradient-flow path; bug~15 is a Bahasa-Malay regex
false-fire in \texttt{priors\_v3} where ``JUMLAH BESAR''
(legitimate \emph{grand-total} synonym) registers as a
distractor; bug~16 is a divide-by-zero in the KD pooling step on
a 0-box receipt; bug~17 is a self-hit leak in the RAG retriever
that artificially inflates val F1 when the query image's own
identifier appears in the retrieval cache.
